A large language model (LLM) pipeline for mapping malware evidence to MITRE ATT\&CK (Adversarial
Tactics, Techniques and Common Knowledge) techniques is mostly not the model. It wraps one in a
sandbox, a disassembler, retrieval indices and rule engines. What the model contributes is therefore
a difference against that tooling, not a score. This study measures that difference component by
component. Three findings are positive. Splitting the analysis across several model calls scores \fact{cape-negotiated-delta} F1 above a
single judge on \fact{cape-consensus-families} real samples \fact{cape-negotiated-ci}, at
\fact{cape-consensus-negotiated-ratio} times its output. A deterministic technique assigner that composes its ranking and gating
backends, rather than choosing between them, gives the cleanest identifier gate on two external
corpora. And turning one decoding flag off gains \fact{reasoning-flag-worth-low} to
\fact{reasoning-flag-worth-high} F1, more than any architecture or parameter count we measured. Four
findings are negative or bounded. The negotiation the design was built for contributes
\fact{cape-mechanism-delta}, at a resolution of \fact{mde-cape-mechanism}: what pays is the extra
calls, not the argument between them. Against the sandbox it is built on, the full pipeline is
\fact{h2h-delta} F1 \fact{h2h-ci}, at a resolution of \fact{h2h-mde}. Three retrieval components move
nothing end to end. And the confidence the model reports, which every deterministic gate consumes,
separates correct claims from wrong ones at an area under the curve of \fact{confidence-auc}, on an
interval containing chance. The paper also reports seven defects in its own instrument, each of
which returned a plausible result rather than an error.
